\documentclass[12pt]{article}
\usepackage{amsmath}
\usepackage{amssymb}
\usepackage[margin=1in]{geometry}

\begin{document}

\begin{center}
\textbf{\Large Homework 4} \\
Tashfeen Hatim Said Omran \\
UCF ID: 5862894 \\
Class: STA4321 \\
November 4, 2026
\end{center}

\vspace{0.5cm}

\noindent \textbf{4.94} \\
\textbf{a)} Here, the random variable $Y \sim \text{Exponential}(\beta = 3.6)$. I need to find the probability that $Y > 9$. 
\[ P(Y > 9) = \int_9^\infty \frac{1}{3.6} e^{-y/3.6} \, dy = \left[ -e^{-y/3.6} \right]_9^\infty = e^{-9/3.6} = e^{-2.5} \approx 0.0821. \]

\vspace{0.2cm}
\noindent \textbf{b)} With the new strategy, $\beta = 2.5$. I need to find the new probability that $Y > 9$.
\[ P(Y > 9) = \int_9^\infty \frac{1}{2.5} e^{-y/2.5} \, dy = \left[ -e^{-y/2.5} \right]_9^\infty = e^{-9/2.5} = e^{-3.6} \approx 0.0273. \]

\vspace{0.5cm}

\noindent \textbf{4.104} \\
Let $p$ be the probability that a single component operates for at least 200 hours. The lifetime is $Y \sim \text{Exponential}(\beta = 100)$.
\[ p = P(Y \ge 200) = \int_{200}^\infty \frac{1}{100} e^{-y/100} \, dy = e^{-200/100} = e^{-2}. \]
The equipment operates if at most 1 component fails out of 3, which means 2 or 3 components must survive. Let $X$ be the number of surviving components. I can model this as a binomial distribution where $X \sim \text{Binomial}(n=3, p=e^{-2})$. I find $P(X \ge 2)$:
\begin{align*}
P(X \ge 2) &= P(X=2) + P(X=3) \\
&= \binom{3}{2}(e^{-2})^2(1-e^{-2})^1 + \binom{3}{3}(e^{-2})^3(1-e^{-2})^0 \\
&= 3e^{-4}(1-e^{-2}) + e^{-6} \\
&= 3e^{-4} - 3e^{-6} + e^{-6} \\
&= 3e^{-4} - 2e^{-6} \approx 0.0500.
\end{align*}

\vspace{0.5cm}

\noindent \textbf{4.131} \\
The measurement error $Y$ has a Beta distribution with $\alpha=1$ and $\beta=2$. I set up the probability density function for $0 < y < 1$:
\[ f(y) = \frac{\Gamma(1+2)}{\Gamma(1)\Gamma(2)} y^{1-1}(1-y)^{2-1} = \frac{2}{1 \cdot 1}(1)(1-y) = 2(1-y). \]
\textbf{a)} I integrate this density function to find $P(Y < 0.5)$:
\[ P(Y < 0.5) = \int_0^{0.5} 2(1-y) \, dy = \left[ 2y - y^2 \right]_0^{0.5} = 2(0.5) - (0.5)^2 = 1 - 0.25 = 0.75. \]

\vspace{0.2cm}
\noindent \textbf{b)} I use the known formulas for the mean and variance of a Beta distribution to find $E(Y)$ and the standard deviation $\sigma$:
\begin{align*}
E(Y) &= \frac{\alpha}{\alpha+\beta} = \frac{1}{1+2} = \frac{1}{3} \approx 0.333 \, \mu\text{s}. \\
V(Y) &= \frac{\alpha\beta}{(\alpha+\beta)^2(\alpha+\beta+1)} = \frac{1(2)}{(1+2)^2(1+2+1)} = \frac{2}{9 \cdot 4} = \frac{2}{36} = \frac{1}{18}. \\
\sigma &= \sqrt{V(Y)} = \sqrt{\frac{1}{18}} = \frac{1}{3\sqrt{2}} \approx 0.2357 \, \mu\text{s}.
\end{align*}

\vspace{0.5cm}

\noindent \textbf{4.136} \\
\textbf{a)} I find the moment-generating function $m(t) = E(e^{tY})$ for $Y \sim \text{Exponential}(\theta)$.
\[ m(t) = \int_0^\infty e^{ty} \left(\frac{1}{\theta}\right) e^{-y/\theta} \, dy = \int_0^\infty \frac{1}{\theta} e^{-y\left(\frac{1}{\theta} - t\right)} \, dy \]
This integral converges for $t < 1/\theta$:
\[ m(t) = \frac{1/\theta}{1/\theta - t} = \frac{1}{1-\theta t} = (1-\theta t)^{-1}. \]

\vspace{0.2cm}
\noindent \textbf{b)} I differentiate $m(t)$ to find $E(Y)$ and $V(Y)$.
\begin{align*}
m'(t) &= -1(1-\theta t)^{-2}(-\theta) = \theta(1-\theta t)^{-2} \implies E(Y) = m'(0) = \theta(1)^{-2} = \theta. \\
m''(t) &= -2\theta(1-\theta t)^{-3}(-\theta) = 2\theta^2(1-\theta t)^{-3} \implies E(Y^2) = m''(0) = 2\theta^2(1)^{-3} = 2\theta^2. \\
V(Y) &= E(Y^2) - [E(Y)]^2 = 2\theta^2 - \theta^2 = \theta^2.
\end{align*}

\vspace{0.5cm}

\noindent \textbf{4.140} \\
I identify the distributions by matching the moment-generating functions to known templates. \\
\textbf{a)} $m(t) = (1-4t)^{-2}$. This matches the Gamma distribution MGF $(1-\beta t)^{-\alpha}$. Thus, the random variable has a \textbf{Gamma distribution with $\alpha = 2$ and $\beta = 4$}. \\
\textbf{b)} $m(t) = 1/(1-3.2t) = (1-3.2t)^{-1}$. This matches the Exponential distribution MGF $(1-\beta t)^{-1}$. Thus, it has an \textbf{Exponential distribution with mean $\beta = 3.2$} (or a Gamma distribution with $\alpha=1, \beta=3.2$). \\
\textbf{c)} $m(t) = e^{-5t+6t^2}$. This matches the Normal distribution MGF $e^{\mu t + \frac{1}{2}\sigma^2t^2}$. I determine $\mu = -5$ and $\frac{1}{2}\sigma^2 = 6$, meaning $\sigma^2 = 12$. Thus, it has a \textbf{Normal distribution with mean $\mu = -5$ and variance $\sigma^2 = 12$}.

\vspace{0.5cm}

\noindent \textbf{4.157} \\
The original component lifetime is $Y \sim \text{Exponential}(100)$, and it is replaced at failure or at 200 hours. This creates a new random variable for time in use, $X = \min(Y, 200)$. \\
\textbf{a)} I define the cumulative distribution function $F_X(x) = P(X \le x)$. \\
For $x < 0$, $F_X(x) = 0$. \\
For $0 \le x < 200$, $X$ takes the value of $Y$, so $P(X \le x) = P(Y \le x) = \int_0^x \frac{1}{100}e^{-y/100} \, dy = 1 - e^{-x/100}$. \\
For $x \ge 200$, the component is strictly replaced by 200 hours, so $X$ never exceeds 200, giving $F_X(x) = 1$.
\[ F_X(x) = \begin{cases} 0 & x < 0 \\ 1 - e^{-x/100} & 0 \le x < 200 \\ 1 & x \ge 200 \end{cases} \]

\vspace{0.2cm}
\noindent \textbf{b)} Because $X$ is a strictly non-negative random variable, I use the tail-integral formula for expected value: $E(X) = \int_0^\infty [1 - F_X(x)] \, dx$.
\begin{align*}
E(X) &= \int_0^{200} (1 - (1 - e^{-x/100})) \, dx + \int_{200}^\infty (1 - 1) \, dx \\
&= \int_0^{200} e^{-x/100} \, dx \\
&= \left[ -100 e^{-x/100} \right]_0^{200} \\
&= -100e^{-2} - (-100e^0) = 100(1 - e^{-2}) \text{ hours}.
\end{align*}

\end{document}